% Part: computability
% Chapter: recursive-functions
% Section: normal-form

\documentclass[../../../include/open-logic-section]{subfiles}

\begin{document}

\olfileid[fr]{cmp}{rec}{nft}
\olsection{Le théorème de forme normale}

\begin{thm}[Théorème de forme normale de Kleene]
\ollabel{thm:kleene-nf}
Il existe une relation récursive primitive $T(e,x,s)$ et une
fonction récursive primitive $U(s)$ ayant la propriété suivante :
pour toute fonction récursive partielle \emph{unaire}~$f$,
il existe un entier~$e$ tel que
\[
f(x) \simeq U(\umin{s}{T(e, x, s)})
\]
pour tout~$x$. Les fonctions de plusieurs arguments peuvent
être traitées en codant leurs tuples d'arguments dans~$x$,
ou au moyen d'une relation~$T$ adaptée à chaque arité.\footnote{Précision éditoriale : la formule source prend un seul argument $x$. Nous explicitons cette arité ; elle représente les autres arités après codage des tuples.}
\end{thm}

\begin{explain}
La démonstration du théorème est complexe, mais son idée
fondamentale est simple. Toute fonction récursive partielle
possède un \emph{indice}~$e$, c'est-à-dire, intuitivement,
un entier qui code son programme ou sa définition. Lorsque
$f(x) \fdefined$, on peut enregistrer systématiquement
son calcul et coder cet enregistrement par un entier~$s$.
Une fonction récursive primitive peut vérifier, à partir de
$x$ et de la définition codée par~$e$, que $s$ est bien
le code d'un calcul terminé. La relation~$T$, qui exprime que
« la fonction d'indice~$e$ possède un calcul sur l'entrée~$x$
et que $s$ code ce calcul », est donc récursive primitive.
Le résultat du calcul, c'est-à-dire la valeur de~$f(x)$,
se déduit lui aussi récursivement primitivement de
l'enregistrement complet~$s$.

Le théorème montre qu'une seule recherche non bornée suffit
à définir toute fonction récursive partielle. Il suffit de
parcourir les entiers jusqu'à trouver le code d'un calcul
de la fonction d'indice~$e$ sur l'entrée~$x$. Nous pouvons
donc employer les entiers~$e$ comme « noms » de fonctions
récursives partielles et noter $\cfind{e}$ la fonction~$f$
définie par l'équation du théorème. Une même fonction peut
avoir plusieurs indices; en fait, toute fonction récursive
partielle en a une infinité dans le codage usuel des programmes, où l'on peut ajouter des instructions inutiles.
\end{explain}

\end{document}
